\subsection{Pengertian Menjaga Lisan}
Menjaga lisan adalah sikap berhati hati dalam berbicara dengan memilih ucapan yang benar, baik, bermanfaat, sopan dan tidak menyakiti orang lain.

\subsection{Dalil Menjaga Lisan}
\begin{enumerate}
\item \textit{Mā yalfizu min qaulin illā ladaihi raqībun ‘atīd.} \\
  \textbf{Artinya: }Tidak ada suatu yang diucapkannya melainkan ada di sisinya malaikat pengawas yang selalu siap mencatat \textbf{Qaf (18)}

\item \textit{Yā ayyuhal-lażīna āmanuttaqullāha wa qūlū qaulan sadīdā.} \\
  \textbf{Artinya: }Wahai orang-orang yang beriman! Bertakwalah kamu kepada Allah dan ucapkanlah perkattan yang benar. \textbf{Al-Ahzab (70)}

\item \textit{‘An Abī Hurairata raḍiyallāhu ‘anhu qāla: qāla Rasūlullāhi ṣallallāhu ‘alaihi wa sallam: Man kāna yu’minu billāhi wal-yaumil-ākhiri falyaqul khairan au liyaṣmut.} \\
  \textbf{Artinya: }Dari Abu Hurairah RA ia berkata Rasulullah SAW bersabda: ``Barangsiapa yang beriman kepada Allah dan hari akhir, hendaklah dia berkata benar atau diam.'' \textbf(HR Muslim \& Bukhari)
\end{enumerate}

\subsubsection{Bentuk-Bentuk Menjaga Lisan}
\begin{enumerate}
\item \textbf{Berkata Jujur} \\
  Tidak mengatakan sesuatu yang bertentangan dengan kenyataan \\
  Contoh: Mengakui kesalahan, tidak membuat alasan palsu atau tidak memalsukan informasi.
\item \textbf{Berkata Baik} \\
  Berkata baik mengandung arti kesopanan, santun, tidak menyakiti dan tidak merendahkan.
\item \textbf{Menghindari Ghibah} \\
  Ghibah adalah membicarakan keburukan seseorang ketika orang tersebut tidak berada di tempat meskipun yang dibicarakan benar.
\item \textbf{Menghindari Fitnah} \\
  Fitnah adalah menuduh/menyebarkan sesuatu yang tidak benar kepada orang lain.
\item \textbf{Menghindari Namimah} \\
  Namimah adalah menyampaikan perkataan seseorang pada orang lain dengan tujuan menimbulkan perselisihan atau permusuhan.
\item \textbf{Menghindari Dusta} \\
  Dusta adalah mengatakan sesuatu yang tidak benar.
\item \textbf{Menghindari Ucapan Kasar} \\
\end{enumerate}

\subsubsection{Bahaya Tidak Menjaga Lisan}
\begin{enumerate}
\item Berdosa.
\item Permusuhan.
\item Hilangnya kepercayaaan.
\item Rusaknya persaudaraan.
\item Rusaknya reputasi seseorang.
\end{enumerate}

\subsubsection{Keutamaan Menjaga Lisan}
\begin{enumerate}
\item Memperoleh pahala.
\item Dipercaya orang lain.
\item Memiliki hubungan sosial yang baik.
\item Memiliki akhlak mulia.
\item Menjaga kehormatan diri dan orang lain.
\end{enumerate}
